\documentclass{article}
\usepackage[a4paper, total={6.5in, 10in}]{geometry}
\setlength{\parindent}{0pt}
\usepackage{tcolorbox}
\tcbset{colback=white!94!black, colframe=black!60!white, coltitle = white, boxrule=0.8pt, arc=2mm}
\newtcolorbox{boxs}[2][]{title= #2,#1}
\usepackage{datetime2}
\usepackage[british]{babel}
\usepackage{amsfonts}
\usepackage{amsmath}

\title{\textbf{Linear Algebra HW 8}}
\author{Robert O'Connell}
\date{\today}
\begin{document}
\maketitle
\vspace{5mm}
\subsection*{Question 1}
\(\chi_A(\lambda)\) is of degree 11 \(\Rightarrow A \in M_{11 \times 11}(\mathbb{C})\)

\begin{itemize}
    \item \(\lambda = 0\), algebraic multiplicity \(6\), largest block associated with is \(3 \times 3\)
    \item \(\lambda = 4\), algebraic multiplicity \(3\), largest block associated with is \(2 \times 2\)
    \item \(\lambda = -5\), algebraic multiplicity \(2\), largest block associated with is \(2 \times 2\)
\end{itemize}

\(\dim(\mathcal{C}(A)) = 8 \Rightarrow \dim(\mathcal{N}(A)) = 3\)

\(\Rightarrow 3\) linearly independent eigenvectors associated with \(\lambda = 0\)

\[
\Rightarrow J = J_2(4) \oplus J_1(4) \oplus J_2(-5) \oplus J_3(0) \oplus J_2(0) \oplus J_1(0)
\]

\subsection*{Question 2}
To check whether the two matrices are similar, I will find the Jordan Normal Form
of each of them and check if it is the same

\(\mathcal{N}(A-I):\)
\[
\begin{bmatrix}
    0 & 1 & 0 & 1 & | 0 \\
    -2 & -1 & 4 & 1 & | 0 \\
    -1 & 0 & 2 & 1 & | 0 \\
    -2 & -3 & 4 & 3 & | 0 \\
\end{bmatrix}
\]
Which row reduces to 
\[
\begin{bmatrix}
    1 & 0 & -2 & -1 & | 0 \\
    0 & 1 & 0 & 1 & | 0 \\
    0 & 0 & 0 & 0 & | 0 \\
    0 & 0 & 0 & 0 & | 0 \\
\end{bmatrix}
\]
Which has two free variables, thus \(\dim(\mathcal{N}(A-I)) = 2\)

\(\mathcal{N}(A-3I):\)
\[
\begin{bmatrix}
    -2 & 1 & 0 & 1 & | 0 \\
    -2 & -3 & 4 & 1 & | 0 \\
    -1 & 0 & 0 & 1 & | 0 \\
    -2 & -3 & 4 & 1 & | 0 \\
\end{bmatrix}
\]
Which row reduces to 
\[
\begin{bmatrix}
    1 & 0 & 0 & -1 & | 0 \\
    0 & 1 & 0 & 3 & | 0 \\
    0 & 0 & 0 & 0 & | 0 \\
    0 & 0 & 0 & 0 & | 0 \\
\end{bmatrix}
\]
Which has two free variables, thus \(\dim(\mathcal{N}(A-3I)) = 2\)

Thus, the Jordan form of \(A\) is
\[
J_A = \begin{bmatrix}
    1 & 0 & 0 & 0 \\
    0 & 1 & 0 & 0 \\
    0 & 0 & 3 & 0 \\
    0 & 0 & 0 & 3 \\
\end{bmatrix}
\]
\\

\(\mathcal{N}(B-I):\)
\[
\begin{bmatrix}
    -2 & 2 & 2 & 0 & | 0 \\
    -2 & 2 & 2 & 0 & | 0 \\
    -2 & 2 & 2 & 0 & | 0 \\
    -2 & 0 & 4 & 2 & | 0 \\
\end{bmatrix}
\]
Which row reduces to 
\[
\begin{bmatrix}
    -2 & 2 & 2 & 0 & | 0 \\
    0 & -2 & 0 & 2 & | 0 \\
    0 & 0 & 0 & 0 & | 0 \\
    0 & 0 & 0 & 0 & | 0 \\
\end{bmatrix}
\]
Which has two free variables, thus \(\dim(\mathcal{N}(B-I)) = 2\)

\(\mathcal{N}(B-3I):\)
\[
\begin{bmatrix}
    -4 & 2 & 2 & 0 & | 0 \\
    -2 & 0 & 2 & 0 & | 0 \\
    -2 & 2 & 0 & 0 & | 0 \\
    -2 & 0 & 4 & 0 & | 0 \\
\end{bmatrix}
\]
Which row reduces to 
\[
\begin{bmatrix}
    -4 & 2 & 2 & 0 & | 0 \\
    0 & -1 & 1 & 0 & | 0 \\
    0 & 0 & 0 & 0 & | 0 \\
    0 & 0 & 0 & 0 & | 0 \\
\end{bmatrix}
\]
Which has two free variables, thus \(\dim(\mathcal{N}(B-3I)) = 2\)

Thus, the Jordan form of \(B\) is
\[
J_B = \begin{bmatrix}
    1 & 0 & 0 & 0 \\
    0 & 1 & 0 & 0 \\
    0 & 0 & 3 & 0 \\
    0 & 0 & 0 & 3 \\
\end{bmatrix}
\]

Since \(A\) and \(B\) have the same Jordan Normal Form, they are similar to each other

\subsection*{Question 3}
If we use the coordinate mapping with respect to the basis \(\mathcal{E} = \{x^3,x^2,x,1\}\), this question becomes one from
\(\mathbb{R} \rightarrow \mathbb{R}\), which is much easier

One such linear operator is then given by the matrix
\[
[T]_{\mathcal{E},\mathcal{E}} = \begin{bmatrix}
    -3 & 0 & 0 & 0 \\
    0 & -3 & 0 & 0 \\
    0 & 0 & 2 & 0 \\
    0 & 0 & 0 & 2 \\
\end{bmatrix}
\]

Then converting back we get
\[
T(ax^3 + bx^2 +cx + d) = -3ax^3 - 3bx^2 + 2cx + 2d
\]

\subsection*{Question 4}
\(\Rightarrow\)

Let \(A\) be an \(n\times n\) invertible matrix, with minimal polynomial \(m(x) = \sum_{i=0}^n a_i x^i\)

\(A\) invertible \(\iff\) \(\dim(\mathcal{N}(A)) = 0\)

\(\iff\) 0 not an eigenvalue of \(A\)

\(\iff\) \(m(0) \neq 0\)

\(\Rightarrow\) \(a_0 \neq 0\)
\\

\(\Leftarrow\)

Let \(A\) be an \(n\times n\) matrix, with minimal polynomial \(m(x) = \sum_{i=0}^n a_i x^i = 
\sum_{i=1}^{n} a_i x^i + a_0\), where \(a_0 \neq 0\)

\(m(A) = 0\) \(\iff\) \(\sum_{i=1}^{n} a_i A^i + a_0I = 0\)

\(\iff\) \(-\frac{1}{a_0}\sum_{i=1}^{n} a_i A^i= I\)

\(\iff\) \(A(-\frac{1}{a_0}\sum_{i=1}^{n} a_i A^{i-1}) = I\)

\(\Rightarrow\) \(AA^{-1} = I\)




\subsection*{Question 5}
\subsubsection*{(a)}
\(A\) is an upper triangular matrix, thus \(\lambda = 1\) is a triple eigenvalue

\(\mathcal{N}(A-I):\)
\[
\begin{bmatrix}
    0 & 1 & -1 & | & 0 \\
    0 & 0 & 0 & | & 0 \\
    0 & 0 & 0 & | & 0 \\
\end{bmatrix}
\]
Which has two free variable, thus \(\dim(\mathcal{N}(A-I)) = 2\)

Thus the Jordan Normal Form of \(A\) is 
\[
J = \begin{bmatrix}
    1 & 1 & 0 \\
    0 & 1 & 0  \\
    0 & 0 & 1  \\
\end{bmatrix}
\]

\subsubsection*{(b)}
\(\chi_A(x) = (1 - x)^3\), thus \(m(x) = \pm (1- x)^k\), for some \(1 \leq k \leq 3\), we also need \(m(A) = 0\)

\[
(I - A) = \begin{bmatrix}
    0 & -1 & 1 \\
    0 & 0 & 0  \\
    0 & 0 & 0  \\    
\end{bmatrix} \neq 0
\]
\[
(I - A)^2 = \begin{bmatrix}
    0 & 0 & 0 \\
    0 & 0 & 0  \\
    0 & 0 & 0  \\
\end{bmatrix}
\]

Thus \(m(x) = (1 - x)^2\)

\subsubsection*{(c)}
\[
A^2 = \begin{bmatrix}
    1 & 2 & -2 \\
    0 & 1 & 0  \\
    0 & 0 & 1  \\
\end{bmatrix}
\]
\[
A^3 = A^2A = \begin{bmatrix}
    1 & 3 & -3 \\
    0 & 1 & 0  \\
    0 & 0 & 1  \\
\end{bmatrix}
\]
Then
\[
A^6 = (A^3)^2 = \begin{bmatrix}
    1 & 6 & -6 \\
    0 & 1 & 0  \\
    0 & 0 & 1  \\
\end{bmatrix}
\]



\subsection*{Question 6}
\subsubsection*{(a)}
\(\mathcal{N}(A - I)\):
\[
\begin{bmatrix}
    -50 & 43 & 7 & | & 0 \\
    -57 & 49 & 8 & | & 0 \\
    -8 & 7 & 1 & | & 0 \\
\end{bmatrix}
\]
Which row reduces to 
\[
\begin{bmatrix}
    1 & 0 & -1 & | & 0 \\
    0 & 1 & -1 & | & 0 \\
    0 & 0 & 0 & | & 0 \\
\end{bmatrix}
\]
Which has one free variable, thus \(\dim(\mathcal{N}(A-I)) = 1\)

Then \(A\) has the Jordan Form
\[
\begin{bmatrix}
    1 & 1 & 0 \\
    0 & 1 & 1 \\
    0 & 0 & 1 \\
\end{bmatrix}
\]

\subsubsection*{(b)}
\(\chi_A(x) = (1 - x)^3\), since a matrix and its Jordan form will have the same characteristic polynomial
thus \(m(x) = \pm (1- x)^k\), for some \(1 \leq k \leq 3\), we also need \(m(A) = 0\)

\[
(I - A) = \begin{bmatrix}
    50 & -43 & -7 \\
    57 & -49 & -8 \\
    8 & -7 & -1 \\ 
\end{bmatrix} \neq 0
\]
\[
(I - A)^2 = \begin{bmatrix}
    -7 & 6 & 1 \\
    -7 & 6 & 1  \\
    -7 & 6 & 1  \\
\end{bmatrix} \neq 0 
\]
\[
(I - A)^3 = \begin{bmatrix}
    0 & 0 & 0 \\
    0 & 0 & 0  \\
    0 & 0 & 0  \\
\end{bmatrix}
\]

Thus \(m(x) = (-1)(1 - x)^3\)

\subsubsection*{(c)}
We know \(m(A)= 0\), thus
\[
\Rightarrow (I - A)^3 = 0
\]
\[
\Rightarrow I - 3A + 3A^2 - A^3 = 0
\]
\[
\Rightarrow A(A^2 - 3A + 3I) = I
\]

\[
\Rightarrow A^{-1} = A^2 - 3A + 3I
\]
\[
\Rightarrow A^{-1} = \begin{bmatrix}
    -106 & 92 & 15 \\
    -121 & 105 & 17 \\
    -23 & 20 & 4
\end{bmatrix} + \begin{bmatrix}
    147 & -129 & -21 \\
    171 & -150 & -24 \\
    24 & -21 & -6
\end{bmatrix} + \begin{bmatrix}
    3 & 0 & 0 \\
    0 & 3 & 0 \\
    0 & 0 & 3
\end{bmatrix}
\]
\[
A^{-1} = \begin{bmatrix}
    142 & -37 & -6 \\
    50 & -42 & -7 \\
    1 & -1 & 1
\end{bmatrix}
\]
\end{document}